% Part:sets-functions-relations
% Chapter: size-of-sets
% Section: introduction

\documentclass[../../../include/open-logic-section]{subfiles}

\begin{document}

\olfileid{sfr}{siz}{int}

\olsection{పరిచయం}

జార్జ్ కాంటర్ 1870లలో సమితి సిద్ధాంతాన్ని అభివృద్ధి చేసినప్పుడు,
అనంత సముదాయం---మధ్యయుగ పండితులు చెప్పినట్లు వాస్తవ అనంతత్వం---అనే
భావనను ఆమోదయోగ్యంగా చేయడం ఆయన లక్ష్యాల్లో ఒకటి. వివిధ సమితుల
\emph{పరిమాణం} గురించి ఆయన ఇచ్చిన విశ్లేషణ ఇందులో కీలక భాగం.
$a$, $b$, $c$ అన్నీ వేర్వేరు అయితే, $\{a, b, c\}$ సమితి సహజబుద్ధికి
$\{a, b\}$ కంటే \emph{పెద్దది}. మరి అనంత సమితుల సంగతేమిటి?
అవన్నీ ఒకదానితో ఒకటి సమాన పరిమాణంలో ఉంటాయా? ఉండవని తేలుతుంది.

ఇక్కడ మొదటి ముఖ్యమైన భావన లెక్కింపు. ప్రతి పరిమిత సమితిలోని
\tetoken{మూలకాలు}{!!{element}s} అన్నింటినీ వరుసగా రాసి జాబితా చేయవచ్చు. జాబితా కూడా
అనంతంగా ఉండనిస్తే, కొన్ని అనంత సమితులలోని \tetoken{మూలకాలు}{!!{element}s} అన్నింటినీ
కూడా జాబితా చేయవచ్చు. అటువంటి సమితులను \tetoken{లెక్కించదగిన}{!!{enumerable}} అంటారు.
ఈ అధ్యాయం చివరికల్లా మనం పూర్తిగా అర్థం చేసుకోబోయే కాంటర్ ఆశ్చర్యకర
ఫలితం ఏమిటంటే, కొన్ని అనంత సమితులు \tetoken{లెక్కించదగిన}{!!{enumerable}} కావు.

\end{document}
